\chapter{Thảo Luận}
\label{chap:discussion}

\section{Diễn Giải Các Phát Hiện Chính}

\subsection{Hiệu Suất Mô Hình và Lựa Chọn Thuật Toán}
Các kết quả thực nghiệm chứng minh rằng tất cả năm thuật toán đã triển khai đạt độ chính xác cao (97,38--99,13\%) trên nhiệm vụ HAR năm lớp, và đạt 100\% trên bài toán 3 lớp cơ bản. Phát hiện này phù hợp với tài liệu gần đây cho thấy rằng các mô hình ML cổ điển được điều chỉnh đúng cách vẫn cạnh tranh với học sâu cho các nhiệm vụ phân loại chuỗi thời gian trên các tập dữ liệu có kích thước vừa phải \cite{anguita2013har, banos2014wisdm}.

Biên độ chênh lệch giữa thuật toán tốt nhất (PyTorch MLP/CNN: 99,13\%) và thấp nhất (SVM: 97,38\%) chỉ 1,75 điểm phần trăm --- cho thấy chất lượng pipeline trích xuất đặc trưng quan trọng hơn lựa chọn thuật toán khi 63 đặc trưng bất biến hướng đã được thiết kế tốt. Huấn luyện nhanh của Random Forest (0,15s) và Neural Network sklearn (0,42s) làm cho chúng hấp dẫn cho tạo mẫu nhanh và các quy trình phát triển lặp lại.

SVM với kernel RBF đạt độ chính xác thấp nhất (97,38\%), với khoảng cách huấn luyện-kiểm tra rất nhỏ (97,55\% vs 97,38\%) gợi ý underfitting thay vì overfitting---kernel RBF với cấu hình mặc định chưa nắm bắt đủ cấu trúc phi tuyến của bài toán 5 lớp. Tối ưu hóa siêu tham số ($C$, $\gamma$) có thể cải thiện đáng kể.

\subsection{Cân Bằng Lớp trong Thực hành}
Cân bằng trọng số lớp (\texttt{class\_weight='balanced'}) được bật mặc định cho RF và SVM trong framework, và \texttt{CrossEntropyLoss} của PyTorch tự xử lý trọng số mẫu. Kết quả thực nghiệm cho thấy các lớp thiểu số (Walking Downstairs: 38 mẫu, Walking Upstairs: 35 mẫu so với Running: 56 mẫu) vẫn đạt recall 0,971--1,000 và F1 $\geq$ 0,949 --- xác nhận hiệu quả của chiến lược cân bằng.

Lựa chọn thiết kế bật cân bằng lớp mặc định phản ánh nguyên tắc ưu tiên triển khai thực tế: trong ứng dụng HAR, tất cả hoạt động phải được phát hiện đáng tin cậy, không chỉ các lớp thường xuyên nhất. Điều này đặc biệt quan trọng trong các bối cảnh an toàn (phát hiện té ngã, giám sát bệnh nhân) nơi bỏ sót lớp thiểu số có thể gây hậu quả nghiêm trọng \cite{he2009learning}.

\subsection{Tính Khả Thi Triển Khai Biên}
Dữ liệu triển khai tham khảo trên bài toán 3 lớp cho thấy mô hình PyTorch 1D-CNN đạt 99,4\% độ chính xác trên thiết bị với tỷ lệ từ chối chỉ 0,6\% (ngưỡng tin cậy $\tau = 0{,}6$) và độ tin cậy trung bình 98,8\%. Kết quả này xác nhận rằng pipeline do framework sinh ra có thể vận hành chính xác trên nền tảng tham chiếu Seeed XIAO nRF52840 (ARM Cortex-M4F @ 64~MHz).

Với cửa sổ 1,5 giây và bước trượt 0,75 giây, phân loại thời gian thực yêu cầu suy luận hoàn thành trong 750ms---biên độ dư thừa đáng kể cho vi điều khiển lớp Cortex-M4F, để lại khoảng trống cho:
\begin{itemize}
    \item Giao tiếp không dây (truyền BLE kết quả)
    \item Quản lý điện năng (ngủ sâu giữa các chu kỳ suy luận)
    \item Các phương thức cảm biến bổ sung
\end{itemize}

Các phép đo chi tiết về thời gian suy luận, bộ nhớ và tiêu thụ điện năng trên nền tảng tham chiếu cần được bổ sung cho phiên bản cuối cùng.

\subsection{Tiền Xử Lý Tương Tác: Một Thay Đổi Mô Hình}
Giao diện cửa sổ thời gian kéo thả mới lạ đại diện cho một thay đổi mô hình từ các quy trình tiền xử lý truyền thống. Các phương pháp thông thường yêu cầu người dùng hoặc:
\begin{enumerate}
    \item Viết mã để phân đoạn dữ liệu theo chương trình (rào cản kỹ thuật cao)
    \item Sử dụng phân đoạn tự động hộp đen (không kiểm soát, lỗi thường xuyên)
    \item Chú thích thủ công các dấu thời gian trong bảng tính (tẻ nhạt, dễ lỗi)
\end{enumerate}

Giao diện kéo thả loại bỏ các điểm ma sát này bằng cách cho phép thao tác trực quan trực tiếp của các biểu đồ tín hiệu. Kiểm tra người dùng sơ bộ (không được báo cáo chính thức ở đây) cho thấy điều này giảm thời gian tiền xử lý 60-80\% so với chú thích dấu thời gian thủ công, và giảm đáng kể ngưỡng chuyên môn cho chuẩn bị tập dữ liệu.

Thiết kế này phù hợp với các nguyên tắc của giao diện thao tác trực tiếp \cite{shneiderman2016designing}, nơi người dùng tương tác với các biểu diễn trực quan của dữ liệu thay vì các đặc tả văn bản trừu tượng. Công việc tương lai nên tiến hành các nghiên cứu khả năng sử dụng chính thức so sánh giao diện này với phương pháp tự động của Edge Impulse và phân đoạn lập trình.

\section{So Sánh với Công Trình Liên Quan}

\subsection{Các Nền Tảng Thương Mại}
So sánh với Edge Impulse và SensiML tiết lộ các đánh đổi giữa chi phí, kiểm soát và tiện lợi:

\textbf{Edge Impulse} ưu tiên dễ sử dụng với kỹ thuật đặc trưng tự động và các quy trình triển khai, nhưng với chi phí đăng ký \$20-100/tháng và tính minh bạch tiền xử lý hạn chế. Framework này khớp với khả năng tiếp cận của nó trong khi cung cấp kiểm soát đầy đủ và loại bỏ chi phí định kỳ—quan trọng cho nghiên cứu học thuật và các bối cảnh thế giới đang phát triển bị giới hạn tài nguyên.

\textbf{SensiML} nhắm mục tiêu triển khai IoT thương mại với khả năng AutoML tinh vi nhưng chi phí \$99--500/tháng và hỗ trợ ít nền tảng hơn.

So sánh số liệu hiệu suất trực tiếp giữa các nền tảng không có ý nghĩa phương pháp luận vì mỗi nền tảng sử dụng tập dữ liệu, pipeline trích xuất đặc trưng và cấu hình thuật toán khác nhau. Thay vào đó, các điểm phân biệt quan trọng nằm ở tính minh bạch, chi phí, và khả năng kiểm soát pipeline (xem Bảng~\ref{tab:qualitative_comparison} trong Chương~\ref{chap:results}).

\subsection{Các Hệ Thống HAR Học Thuật}
So với nghiên cứu HAR học thuật sử dụng các phương pháp dựa trên IMU tương tự \cite{anguita2013har}, độ chính xác đạt được (97--99\% trên 5 lớp) cạnh tranh với các kết quả đã công bố trên các tập dữ liệu benchmark (92--97\% điển hình). Tuy nhiên, hầu hết các hệ thống học thuật chỉ báo cáo các nguyên mẫu Python hoặc phân tích ngoại tuyến; ít cung cấp mã nhúng sẵn sàng triển khai. Đóng góp của luận văn nằm ở việc bắc cầu ``khoảng cách triển khai'' này bằng cách tự động hóa việc dịch từ các mô hình scikit-learn/PyTorch đã huấn luyện sang mã suy luận vi điều khiển.

Cần lưu ý rằng kết quả 97--99\% được đo trên dữ liệu đối tượng đơn; hiệu suất trên tập dữ liệu đa đối tượng (như UCI-HAR) có thể thấp hơn do biến thiên giữa các cá nhân.

\section{Các Hạn Chế và Mối Đe Dọa đối với Tính Hợp Lệ}

\subsection{Hạn Chế Tập Dữ Liệu}
Thực nghiệm xác thực sử dụng dữ liệu từ một đối tượng đơn thực hiện năm hoạt động. Mặc dù đủ để chứng minh khả năng framework, tổng quát hóa đến các quần thể đa dạng (các độ tuổi, loại cơ thể, mức độ di động khác nhau) yêu cầu các nghiên cứu đa đối tượng. Các lớp hoạt động cũng hạn chế—các ứng dụng thực tế có thể cần phát hiện:
\begin{itemize}
    \item Sự kiện té ngã (quan trọng cho chăm sóc người cao tuổi)
    \item Các hoạt động phức tạp (nấu ăn, dọn dẹp, lái xe)
    \item Chuyển đổi giữa các hoạt động
    \item Các mẫu dáng đi bất thường (chấn thương, tình trạng thần kinh)
\end{itemize}

Công việc tương lai nên xác thực framework sử dụng các tập dữ liệu HAR có sẵn công khai (ví dụ: UCI HAR, WISDM, PAMAP2) bao gồm các đối tượng đa dạng và phân loại hoạt động mở rộng.

\subsection{Lựa Chọn Kích Thước Cửa Sổ}
Cửa sổ 1,5 giây (150 mẫu @ 100Hz) với bước trượt 0,75 giây là một lựa chọn cân bằng bối cảnh thời gian và khả năng phản hồi. Nghiên cứu trước đây cho thấy kích thước cửa sổ ảnh hưởng đáng kể đến độ chính xác \cite{banos2014wisdm}: các cửa sổ ngắn hơn (<0.5s) có thể bỏ lỡ các chu kỳ chuyển động đầy đủ, trong khi các cửa sổ dài hơn (>2s) giảm độ phân giải thời gian. Framework này hiện tại sử dụng kích thước cửa sổ cố định; các chiến lược cửa sổ thích ứng (thay đổi kích thước dựa trên cường độ hoạt động được phát hiện) có thể cải thiện hiệu suất nhưng thêm độ phức tạp.

\subsection{Vị Trí và Hướng Cảm Biến}
Các thử nghiệm giả định IMU được đeo ở cổ tay trong một hướng nhất quán. Triển khai thực tế gặp phải sự biến đổi:
\begin{itemize}
    \item Các vị trí cơ thể khác nhau (cổ tay, hông, mắt cá chân) tạo ra các đặc tính tín hiệu riêng biệt
    \item Quay thiết bị tương đối với các trục cơ thể giới thiệu sự mơ hồ về hướng
    \item Nhiều cảm biến đồng thời (ví dụ: cả hai cổ tay + hông) yêu cầu hợp nhất cảm biến
\end{itemize}

Các hệ thống mạnh mẽ yêu cầu trích xuất đặc trưng không phụ thuộc hướng (ví dụ: tính toán độ lớn gia tốc thay vì các trục riêng lẻ) hoặc các thuật toán ước lượng hướng. Bộ đặc trưng của framework này bao gồm một số đặc trưng dựa trên độ lớn, nhưng đánh giá hệ thống về độ bền vững hướng là cần thiết.

\subsection{Sự Đa Dạng Nền Tảng Tính Toán}
Framework hiện đã hiện thực generator cho nhiều họ đích, nhưng mức độ xác thực cần được phân tách thành hai lớp. Ở lớp thứ nhất, \textit{khả năng triển khai} đã được chứng minh bởi kiến trúc factory, các generator đặc thù nền tảng và các gói mã đầu ra tương ứng. Ở lớp thứ hai, \textit{benchmark định lượng liên thiết bị} vẫn mới chỉ được đo chi tiết trên XIAO. Vì vậy, phát biểu chính xác của luận văn không nên là ``framework đã benchmark trên mọi thiết bị'', mà là ``framework cung cấp năng lực triển khai đa thiết bị ở mức kiến trúc và sinh mã, với XIAO là nền tảng thực nghiệm tham chiếu đầu tiên''. Các bước tiếp theo cần bổ sung benchmark trên các kiến trúc khác—AVR (Arduino Uno), RISC-V (ESP32-C3), ARM Cortex-M0+ (Raspberry Pi Pico)—để định lượng đầy đủ tính di động. Đặc điểm bộ nhớ và tốc độ chắc chắn sẽ khác nhau, đặc biệt trên các bộ xử lý AVR 8-bit thiếu hỗ trợ dấu phẩy động phần cứng.

\subsection{Tính Hợp Lệ Bên Ngoài}
Kết quả cụ thể cho nhận dạng hoạt động con người sử dụng các cảm biến IMU đeo cổ tay. Khả năng áp dụng cho các lĩnh vực theo dõi chuyển động khác (nhận dạng cử chỉ, theo dõi hành vi động vật, lập hồ sơ chuyển động robot) vẫn cần được xác thực. Kiến trúc của framework được thiết kế để không phụ thuộc miền, nhưng các chiến lược tiền xử lý và bộ đặc trưng có thể yêu cầu thích ứng cho các ứng dụng không phải HAR.

\section{Các Đánh Đổi Thiết Kế và Các Lựa Chọn Thay Thế}

\subsection{ML Cổ Điển vs Học Sâu}
Luận văn này cố ý chọn học máy cổ điển (Random Forest, SVM, Neural Networks) thay vì học sâu hiện đại (CNN, LSTM, Transformer) do:
\begin{itemize}
    \item \textbf{Ràng Buộc Triển Khai}: Các mô hình cổ điển biên dịch thành mã C++ nhỏ, trong khi các mạng sâu yêu cầu các runtime suy luận chuyên biệt (TFLite Micro) thêm 150--250KB overhead
    \item \textbf{Hiệu Quả Huấn Luyện}: Các mô hình cổ điển huấn luyện trong dưới 1 giây trên CPU; PyTorch huấn luyện trong 2--10 giây
    \item \textbf{Khả Năng Diễn Giải}: Tầm quan trọng đặc trưng Random Forest và ranh giới quyết định SVM có thể kiểm toán; các mạng sâu mờ đục
\end{itemize}

Tuy nhiên, kết quả thực nghiệm cho thấy PyTorch MLP và 1D-CNN đạt hiệu suất cao nhất (99,13\% trên 5 lớp), xác nhận giá trị của cả hai hướng tiếp cận. Framework hỗ trợ đồng thời cả hai: ML cổ điển cho prototyping nhanh và vi điều khiển hạn chế, học sâu cho bài toán phức tạp hơn.

\subsection{Đặc Trưng Thủ Công vs Đặc Trưng Học}
Bộ 63 đặc trưng bất biến hướng (orientation-invariant) được sử dụng bao gồm thống kê thời gian trên biên độ tổng hợp (acc\_mag, gyro\_mag, acc\_jerk\_mag) cộng với đặc trưng phổ DFT. Thiết kế này yêu cầu kiến thức miền nhưng đảm bảo khả năng diễn giải, biểu diễn gọn nhẹ, bất biến hướng gắn thiết bị, và tương đồng bit-exact giữa Python và C++. Kết quả thực nghiệm (97--99\% trên 5 lớp) xác nhận hiệu quả của bộ đặc trưng này. Học sâu từ đầu đến cuối (1D-CNN) đạt hiệu suất tương đương (99,13\%) nhưng với chi phí mô hình lớn hơn và tính minh bạch giảm.

Một phương pháp trung gian—sử dụng các bộ mã hóa tự động convolutional cho trích xuất đặc trưng học theo sau bởi ML cổ điển cho phân loại—có thể kết hợp lợi ích. Điều này vẫn là công việc tương lai.

\subsection{Tạo Mã Trực tiếp vs Runtime Suy luận: Đánh Giá Thực tiễn}
\label{sec:codegen_tradeoff_discussion}

Phần~\ref{sec:codegen_rationale} trình bày lý do thiết kế cho việc chọn tạo mã trực tiếp thay vì runtime suy luận. Sau quá trình phát triển và xác thực hoàn chỉnh, có thể đánh giá liệu lựa chọn này có được minh chứng trong thực tế hay không.

\paragraph{Lợi ích Đã được Xác nhận}
Ba lợi ích dự kiến đã được xác nhận qua thực nghiệm:

Thứ nhất, \textit{dấu chân bộ nhớ tối thiểu}: Neural Network chỉ chiếm 95KB flash---bao gồm trọng số mô hình, logic suy luận, trích xuất đặc trưng 15 thống kê, và chuẩn hóa StandardScaler. Với TFLite Micro, cùng chức năng sẽ yêu cầu thêm 150--250KB cho runtime interpreter, đẩy tổng lên 245--345KB và chiếm $>$25\% flash khả dụng chỉ cho suy luận. Trên nền tảng tham chiếu Seeed XIAO (1MB flash), điều này vẫn khả thi; trên các vi điều khiển với 256KB flash (ví dụ: STM32F103, ATmega2560), phương pháp runtime sẽ không thể triển khai.

Thứ hai, \textit{minh bạch cho gỡ lỗi}---đây là yếu tố có tác động lớn nhất ngoài dự kiến. Khi mô hình đã triển khai luôn dự đoán \texttt{walking\_downstairs}, có thể: (a) đọc chính xác mã C++ đã tạo để xác minh logic suy luận, (b) in giá trị đặc trưng trung gian trên serial để so sánh với Python, (c) xác định lỗi zero-padding và kurtosis formula bằng cách truy vết từng dòng. Với runtime đóng gói, bước (b) và (c) sẽ yêu cầu công cụ debugging chuyên biệt hoặc không thể thực hiện. Kinh nghiệm phát triển thực tế cho thấy 3 trong 10 phát hiện kỹ thuật (Finding 1: zero-padding, Finding 2: kurtosis, Finding 7: feature order) chỉ có thể phát hiện nhờ khả năng kiểm tra mã đã tạo dòng-theo-dòng.

Thứ ba, \textit{hỗ trợ trực tiếp ML cổ điển}: Framework hỗ trợ Random Forest---thuật toán không có đường dẫn chuyển đổi trực tiếp sang TFLite Micro. Chuỗi scikit-learn $\rightarrow$ ONNX $\rightarrow$ TFLite giới thiệu rủi ro sai lệch; trong khi phương pháp trực tiếp đọc \texttt{model.estimators\_} (danh sách cây quyết định) và tạo mã duyệt cây nhị phân chính xác.

\paragraph{Chi phí Phát triển Thực tế}
Tạo mã trực tiếp cũng có chi phí không nhỏ. Framework yêu cầu triển khai và xác minh logic suy luận riêng cho mỗi thuật toán: lan truyền tiến cho Neural Network, tính kernel RBF và bỏ phiếu OvO cho SVM, duyệt cây nhị phân và bỏ phiếu đa số cho Random Forest. Tổng cộng, \texttt{base\_generator.py} chứa $\sim$1580 dòng mã cho trích xuất đặc trưng chung, và mỗi trình tạo đặc thù thêm 200--500 dòng. Bất kỳ lỗi nào trong logic suy luận đã tạo đều ảnh hưởng trực tiếp đến độ chính xác triển khai---như đã chứng minh bởi Finding 2 (kurtosis) và Finding 7 (feature order).

Tuy nhiên, chi phí này là \textit{một lần}: sau khi trình tạo được xác minh cho một thuật toán, nó hoạt động chính xác cho mọi mô hình được huấn luyện bằng thuật toán đó, bất kể số lớp, số đặc trưng, hoặc siêu tham số.

\paragraph{Kết luận về Đánh đổi}
Tạo mã trực tiếp là lựa chọn đúng cho \textit{bối cảnh sử dụng cụ thể của framework}: tập dữ liệu vừa phải ($<$10K mẫu), thuật toán ML cổ điển, vi điều khiển bị giới hạn tài nguyên, và yêu cầu minh bạch cao. Đối với các bối cảnh khác---tập dữ liệu lớn, mạng sâu, vi điều khiển $>$4MB flash---runtime suy luận (TFLite Micro) hoặc trình biên dịch mô hình (TVM) sẽ phù hợp hơn. Framework được thiết kế modular để tích hợp đường dẫn TFLite Micro trong tương lai mà không thay thế phương pháp trực tiếp hiện có.

\subsection{Các Chế Độ Tối Ưu Hóa}
Bốn chế độ tối ưu hóa (accuracy, speed, power, balanced) cung cấp tính linh hoạt, nhưng yêu cầu người dùng hiểu các đánh đổi. Các thiết kế thay thế có thể:
\begin{itemize}
    \item \textbf{Lập Hồ Sơ Tự Động}: Framework đo khả năng thiết bị và chọn chế độ tối ưu
    \item \textbf{Tối Ưu Hóa Đa Mục Tiêu}: Khám phá biên Pareto để tìm các thỏa hiệp độ chính xác-tốc độ-điện năng tốt nhất
    \item \textbf{Runtime Thích Ứng}: Thiết bị điều chỉnh động độ phức tạp suy luận dựa trên mức pin hoặc cường độ chuyển động được phát hiện
\end{itemize}

Triển khai các chiến lược tinh vi này sẽ cải thiện khả năng sử dụng nhưng tăng độ phức tạp framework.

\section{Tương đồng Huấn luyện-Triển khai trong Edge ML}
\label{sec:training_deployment_parity}

Một phát hiện quan trọng trong quá trình phát triển và xác thực framework là tầm quan trọng thiết yếu của \textit{tương đồng huấn luyện-triển khai} (training-deployment parity)---sự đảm bảo rằng quá trình trích xuất đặc trưng trong môi trường huấn luyện (Python) tạo ra kết quả \textbf{giống hệt} với quá trình trích xuất đặc trưng trên thiết bị biên (C++/MicroPython). Vấn đề này ít được nghiên cứu trong tài liệu edge ML nhưng có tác động nghiêm trọng đến hiệu suất triển khai thực tế \cite{paleyes2022challenges, sculley2015hidden}.

\subsection{Các Nguồn Không khớp Được Xác định}
Trong quá trình phát triển, đã xác định và giải quyết nhiều nguồn sai lệch giữa pipeline huấn luyện Python và mã suy luận C++ đã tạo:

\paragraph{Artifact Đệm Số Không (Zero-Padding Artifact)}
Pipeline trích xuất đặc trưng ban đầu đệm các cửa sổ ngắn hơn kích thước mục tiêu bằng hàng toàn số không. Trên thiết bị thực, bộ đệm trượt luôn chứa đầy dữ liệu cảm biến thực (không có giá trị zero), vì cảm biến liên tục đo gia tốc trọng trường. Sự khác biệt này tạo ra phân phối đặc trưng hoàn toàn khác nhau giữa huấn luyện ($\text{acc\_mag\_min} = 0$) và triển khai ($\text{acc\_mag\_min} \approx 9{,}81$), đẩy tất cả đặc trưng đầu vào ra ngoài vùng huấn luyện.

\paragraph{Sai lệch Công thức Thống kê}
Công thức độ nhọn (kurtosis) trong mã C++ tạo ra ban đầu sử dụng độ lệch chuẩn mẫu ($n-1$ ở mẫu số) cho chuẩn hóa z-scores, trong khi Python pandas sử dụng độ lệch chuẩn tổng thể ($n$ ở mẫu số). Sai lệch hệ thống:
\begin{equation}
    \text{Ratio} = \left(\frac{n-1}{n}\right)^2 \approx 0{,}987 \text{ cho } n = 150
\end{equation}
tương đương sai lệch $\sim$1,3\% cho mỗi đặc trưng kurtosis và skewness---nhỏ nhưng hệ thống và tích lũy qua 6 trục.

\paragraph{Thứ tự Đặc trưng Không khớp}
Python huấn luyện trên các cột DataFrame được sắp xếp theo thứ tự bảng chữ cái, trong khi C++ trích xuất đặc trưng theo thứ tự cố định (magnitude $\rightarrow$ per-axis). Mức reorder đặc trưng tại thời điểm tạo mã đảm bảo trọng số mô hình và tham số scaler tương ứng chính xác với thứ tự trích xuất C++.

\subsection{Tác động Thực tế}
Trong ca kiểm chứng trên thiết bị tham chiếu Seeed XIAO nRF52840, các lỗi không khớp khiến mô hình luôn dự đoán lớp \texttt{walking\_downstairs} bất kể hoạt động thực tế. Phân tích cho thấy output biases của mạng nơ-ron $[-0{,}116; -0{,}160; -0{,}066; 0{,}116; 0{,}048]$ chi phối khi đặc trưng đầu vào nằm hoàn toàn ngoài phân phối huấn luyện---lớp có bias cao nhất ($0{,}116$, tương ứng \texttt{walking\_downstairs}) luôn thắng quyết. Bài học này không chỉ áp dụng cho XIAO mà cho toàn bộ ma trận đích của framework, vì bất kỳ backend nào cũng sẽ thất bại nếu tính tương đồng huấn luyện-triển khai không được đảm bảo.

\subsection{Danh sách Kiểm tra Tương đồng}
Bảng~\ref{tab:parity_checklist} tóm tắt 12 biện pháp đảm bảo tương đồng đã được triển khai và xác minh trong framework:

\begin{table}[H]
    \centering
    \caption{Danh sách kiểm tra tương đồng huấn luyện-triển khai}
    \small
    \begin{tabular}{@{}clp{7cm}@{}}
        \toprule
        \# & Biện pháp & Mô tả \\
        \midrule
        1 & Loại bỏ FFT & Tránh sai lệch triển khai FFT giữa Python và C++ \\
        2 & Edge-value replication & Thay thế zero-padding bằng lặp giá trị biên \\
        3 & Kurtosis/skewness & Sử dụng population std cho z-scores, khớp pandas \\
        4 & Thứ tự đặc trưng & Reorder tham số mô hình tại thời điểm tạo mã \\
        5 & Đơn vị cảm biến & m/s² cho gia tốc, deg/s cho con quay hồi chuyển \\
        6 & Tốc độ lấy mẫu & 100Hz trong cả thu thập và suy luận \\
        7 & Kích thước cửa sổ & 150 mẫu, đọc từ metadata mô hình \\
        8 & StandardScaler & Cùng mean/std, cùng clamp range \\
        9 & Chuẩn hóa đơn & Scaling chỉ xảy ra 1 lần trong pipeline \\
        10 & Feature order remap & Trọng số/scaler reorder cho thứ tự C++ \\
        11 & Trích xuất idempotent & Tham số mô hình chỉ trích xuất 1 lần \\
        12 & Xác thực kiến trúc & Phân biệt CNN (raw window) vs feature-based \\
        \bottomrule
    \end{tabular}
    \label{tab:parity_checklist}
\end{table}

\subsection{So sánh với Nền tảng Thương mại}
Các nền tảng thương mại như Edge Impulse hoạt động theo mô hình hộp đen (black-box), trong đó pipeline trích xuất đặc trưng không thể kiểm tra hoặc xác minh bởi người dùng. Nếu một lỗi tương đồng tương tự tồn tại, người dùng không có cách nào phát hiện hoặc sửa chữa. Tính \textbf{minh bạch hoàn toàn} (full transparency) của framework cho phép:
\begin{itemize}
    \item Kiểm tra trực tiếp giá trị đặc trưng (phát hiện $\text{acc\_mag\_min} = 0$)
    \item Truy vết ngược đến nguyên nhân gốc (zero-padding)
    \item Xác minh từng công thức giữa Python và C++
    \item Sửa lỗi và huấn luyện lại mà không chờ nhà cung cấp
\end{itemize}

Phát hiện này nhấn mạnh rằng \textbf{tính minh bạch không chỉ là triết lý mã nguồn mở mà là yêu cầu kỹ thuật thiết yếu} cho triển khai edge ML đáng tin cậy---đặc biệt khi lỗi biểu hiện ở cấp thiết bị chứ không phải ở cấp mô hình.

\section{Tăng cường Dữ liệu và Thiết kế Nhận biết Lớp}
\label{sec:augmentation_discussion}

\subsection{Hiệu quả Tăng cường cho Dữ liệu IMU}
Tăng cường dữ liệu là kỹ thuật đã được chứng minh hiệu quả trong cải thiện tổng quát hóa mô hình, đặc biệt với tập dữ liệu nhỏ \cite{um2017data, iwana2021empirical}. Tuy nhiên, việc áp dụng cho dữ liệu IMU đòi hỏi cân nhắc cẩn thận vì tín hiệu cảm biến quán tính có các ràng buộc vật lý nghiêm ngặt: gia tốc trọng trường luôn hiện diện, biên độ dao động phải nằm trong phạm vi vật lý, và tỷ lệ tín hiệu-trên-nhiễu (SNR) giữa các hoạt động khác nhau rất đáng kể.

Năm phương pháp tăng cường được chọn mô phỏng các nguồn biến thiên thực tế mà mô hình sẽ gặp trong triển khai: nhiễu cảm biến (jitter), biến thiên cường độ (scaling), thay đổi vị trí đặt cảm biến (rotation), biến thiên tốc độ thực hiện (time warping) và tính bất biến với thứ tự thời gian (permutation). Sự kết hợp này bao phủ các nguồn biến thiên chính được xác định trong tài liệu HAR \cite{eyobu2018feature, rashid2022ahar}.

\subsection{Bài học từ Tăng cường Không-Nhận-biết Lớp}
Thực nghiệm ban đầu áp dụng tất cả năm phương pháp đồng nhất cho mọi lớp hoạt động, phát hiện vấn đề: mô hình Neural Network sau tăng cường thống nhất phân loại sai hoạt động \textit{still} thành \textit{walking\_downstairs} với tỷ lệ đáng kể. Nguyên nhân gốc: jitter ($\sigma = 0{,}05$) và rotation ($15^{\circ}$) trên cửa sổ tĩnh tạo ra dao động nhân tạo có biên độ tương đương hoạt động cường độ thấp (đi xuống cầu thang chậm), dẫn đến chồng lấn vùng đặc trưng giữa hai lớp.

Phân tích định lượng xác nhận: perturbation tối đa trên cửa sổ tĩnh với tăng cường đầy đủ đạt 0,6059 (tương đương dao động gia tốc đáng kể), trong khi cửa sổ tĩnh gốc có perturbation $< 0{,}01$. Chiến lược nhận biết lớp giảm perturbation tĩnh về $\leq 0{,}0002$ (micro-jitter), bảo toàn hoàn toàn đặc tính phân biệt.

\subsection{Nguyên tắc Thiết kế}
Kinh nghiệm này dẫn đến nguyên tắc thiết kế tăng cường cho IMU HAR:
\begin{enumerate}
    \item \textbf{Bảo toàn đặc tính phân biệt}: Tăng cường không được phá hủy đặc tính khiến một lớp khác biệt với các lớp khác---với hoạt động tĩnh, đó là biên độ dao động cực thấp
    \item \textbf{Phân biệt theo tính chất hoạt động}: Hoạt động động và tĩnh có bản chất vật lý khác nhau cơ bản; áp dụng cùng phép biến đổi là không phù hợp
    \item \textbf{Phát hiện tự động}: Framework nên phát hiện tự động lớp tĩnh/động thay vì yêu cầu người dùng cấu hình thủ công, giảm rào cản kỹ thuật
\end{enumerate}

Quan điểm này khác biệt với đa số nghiên cứu tăng cường dữ liệu cho HAR, vốn áp dụng biến đổi đồng nhất \cite{um2017data, le_guennec2016augmentation}. Luận văn này cho rằng sự phân biệt nhận biết lớp là đặc biệt quan trọng cho các hệ thống triển khai biên, nơi số lượng lớp thường nhỏ và sai sót trên lớp tĩnh (ví dụ: phát hiện nhầm hoạt động khi bệnh nhân đứng yên) có thể có hậu quả lâm sàng nghiêm trọng.

\section{Từ chối Hoạt động Dựa trên Ngưỡng Tin cậy}
\label{sec:confidence_discussion}

\subsection{Vấn đề Phân loại Ép buộc}
Các hệ thống phân loại tiêu chuẩn gán một nhãn cho mọi đầu vào, kể cả khi đầu vào thuộc phân phối ngoài (out-of-distribution). Trong triển khai HAR thực tế, đây là vấn đề phổ biến: người dùng thường thực hiện các hoạt động không thuộc tập huấn luyện (ngồi xuống ghế, gãi đầu, cúi lấy đồ), và cảm biến ghi nhận tín hiệu chuyển động tương ứng. Bộ phân loại không có cơ chế từ chối sẽ phân loại sai các đầu vào này với độ tin cậy biểu kiến cao---hiện tượng được gọi là \textit{overconfident extrapolation} \cite{hendrycks2017baseline}.

\subsection{Phương pháp Tin cậy Trực tiếp}
Thay vì huấn luyện một mô hình phát hiện bất thường riêng biệt (approach của Edge Impulse), framework trích xuất tin cậy trực tiếp từ đầu ra phân loại hiện có. Phương pháp này có ba ưu điểm thực tế:
\begin{itemize}
    \item \textbf{Không tốn thêm bộ nhớ}: Không mô hình bổ sung, không tham số bổ sung
    \item \textbf{Không cần dữ liệu bổ sung}: Không yêu cầu thu thập mẫu ``không hoạt động'' để huấn luyện
    \item \textbf{Tích hợp tự nhiên}: Tin cậy được tính trong quá trình suy luận đã có, không thêm bước xử lý
\end{itemize}

\subsection{Hạn chế và Hướng Tương lai}
Cần lưu ý rằng tin cậy softmax thường \textit{kém hiệu chỉnh} (poorly calibrated) \cite{guo2017calibration}---giá trị softmax 0,8 không nhất thiết có nghĩa 80\% xác suất đúng. Điều này có nghĩa:
\begin{itemize}
    \item Random Forest có tin cậy hiệu chỉnh tự nhiên tốt hơn (tỷ lệ bỏ phiếu có ý nghĩa thống kê xác suất)
    \item Neural Network/SVM qua softmax có thể quá tự tin---các kỹ thuật hiệu chỉnh nhiệt độ (temperature scaling) \cite{guo2017calibration} có thể cải thiện chất lượng tin cậy
    \item Ngưỡng $\tau = 0{,}6$ là giá trị thực nghiệm; điều chỉnh tự động dựa trên đường cong ROC tập xác thực là hướng tương lai
\end{itemize}

Dù vậy, ngay cả tin cậy chưa hiệu chỉnh cũng cung cấp ``tuyến phòng thủ đầu tiên'' hiệu quả chống lại các dự đoán sai nghiêm trọng, đặc biệt cho các đầu vào hoàn toàn ngoài phân phối nơi tất cả xác suất lớp đều thấp.

\section{Ý Nghĩa cho Thực Hành}

\subsection{Nghiên Cứu và Giáo Dục}
Framework dân chủ hóa nghiên cứu edge AI bằng cách loại bỏ rào cản chi phí (so với các nền tảng thương mại \$20-500/tháng) và rào cản kỹ thuật (so với triển khai TensorFlow Lite thủ công). Điều này cho phép:
\begin{itemize}
    \item Các khóa học ML đại học/sau đại học bao gồm các phòng thí nghiệm triển khai biên thực hành
    \item Các nhà nghiên cứu ở các nước đang phát triển tiến hành các nghiên cứu HAR mà không cần ngân sách thể chế cho giấy phép thương mại
    \item Tạo mẫu nhanh các thuật toán nhận dạng hoạt động mới lạ mà không cần lập trình nhúng cấp thấp
\end{itemize}

\subsection{Các Ứng Dụng Y Tế}
Các mô hình đã xác thực có thể hỗ trợ giám sát sức khỏe đeo được:
\begin{itemize}
    \item \textbf{Phục Hồi Chức Năng}: Theo dõi tiến trình phục hồi bằng cách giám sát mức độ hoạt động và chất lượng dáng đi
    \item \textbf{Chăm Sóc Người Cao Tuổi}: Phát hiện té ngã, hành vi ít vận động và suy giảm hoạt động
    \item \textbf{Thử Nghiệm Lâm Sàng}: Đo lường hoạt động khách quan thay thế nhật ký tự báo cáo
\end{itemize}

Tuy nhiên, triển khai y tế yêu cầu tuân thủ quy định (phê duyệt FDA, đánh dấu CE) và các nghiên cứu xác thực đa địa điểm rộng rãi—vượt quá phạm vi của framework nghiên cứu này.

\subsection{Thể Thao và Thể Dục}
Các thiết bị đeo tiêu dùng (đồng hồ thông minh, băng thể dục) có thể tích hợp các mô hình được tạo bởi framework cho:
\begin{itemize}
    \item Phát hiện tập luyện tự động (chạy, đạp xe, cử tạ)
    \item Phân tích hình thức (phát hiện kỹ thuật squat không đúng, bất đối xứng dáng chạy)
    \item Huấn luyện cá nhân hóa (phản hồi thời gian thực về cường độ hoạt động)
\end{itemize}

\section{Phân Tích Các Phát Hiện Kỹ Thuật Quan Trọng}

\subsection{Bộ lọc Kalman: Đột Phá về Training-Deployment Parity}
\label{sec:disc_kalman_parity}

Việc triển khai bộ lọc Kalman constant-velocity đại diện cho một đột phá quan trọng trong việc giải quyết vấn đề khoảng cách tương đồng huấn luyện-triển khai (training-deployment parity gap). Đây là lần đầu tiên framework có một bộ lọc tiền xử lý với \textbf{khoảng cách parity bằng không}.

\subsubsection{So sánh với Các Phương pháp IIR Truyền thống}
Các bộ lọc IIR Butterworth truyền thống tồn tại lỗ hổng parity vốn có:
\begin{itemize}
    \item \textbf{Huấn luyện}: \texttt{filtfilt} (non-causal, zero-phase) xử lý dữ liệu theo hai chiều
    \item \textbf{Triển khai}: \texttt{lfilter} (causal, forward-only) chỉ có thể xử lý theo một chiều do ràng buộc thời gian thực
\end{itemize}

Khoảng cách này gây sai lệch đặc biệt nghiêm trọng với các đặc trưng nhạy cảm với pha như skewness và kurtosis, có thể dẫn đến suy giảm hiệu suất mô hình 2-5\% khi triển khai thực tế.

\subsubsection{Ưu thế Thiết kế của Kalman Filter}
Bộ lọc Kalman khắc phục vấn đề này từ gốc thông qua thiết kế vốn có:
\begin{enumerate}
    \item \textbf{Tính nhân quả tự nhiên}: Kalman filter về bản chất là causal---chỉ sử dụng observation hiện tại và trạng thái trước đó
    \item \textbf{Tính thích ứng}: Gain điều chỉnh tự động dựa trên prediction error, tạo ra đặc tính lọc thích ứng
    \item \textbf{Mô hình vật lý}: Constant-velocity model phù hợp với bản chất chuyển động của dữ liệu IMU
    \item \textbf{Tham số trực quan}: Process noise (Q) và measurement noise (R) có ý nghĩa vật lý rõ ràng
\end{enumerate}

\textbf{So sánh với Edge Impulse}: Edge Impulse không cung cấp Kalman filtering như một tùy chọn tiền xử lý. Spectral analysis block của họ sử dụng các bandpass filter cố định, thiếu khả năng thích ứng và đảm bảo parity của Kalman filter.

\subsection{Phân Tích Kiến trúc Tạo Mã Đa Kiến trúc}
\label{sec:disc_multi_arch}

Việc hỗ trợ đồng thời các mô hình feature-based (RF/SVM/NN) và end-to-end (CNN) trong cùng một framework đặt ra những thách thức kiến trúc độc đáo mà Finding 16 đã chỉ ra.

\subsubsection{Sự Phức Tạp Metadata theo Kiến trúc}
Case study về CNN metadata mismatch (Finding 16) minh họa tầm quan trọng của architecture-aware code generation:

\textbf{Vấn đề}: CNN model hiển thị filename "f33" (33 features) mặc dù thực tế sử dụng 6 channels per timestep.

\textbf{Nguyên nhân gốc}: Fallback logic không phân biệt được giữa:
\begin{itemize}
    \item Feature-based models: cần statistical feature names (e.g., "acc\_x\_mean", "gyro\_y\_kurtosis")  
    \item CNN models: cần channel placeholders (e.g., "ch0", "ch1", ..., "ch5")
\end{itemize}

\textbf{Giải pháp kiến trúc}: Architecture-aware metadata generation với logic phân nhánh rõ ràng cho từng loại mô hình.

\subsubsection{Phân Biệt với Edge Impulse}
Edge Impulse xử lý sự phức tạp này bằng cách có "processing blocks" riêng biệt cho từng loại mô hình. Framework này phải quản lý cả hai pipeline trong cùng một codebase thống nhất, đòi hỏi:
\begin{itemize}
    \item Validation logic nhận biết kiến trúc
    \item Metadata semantics khác nhau cho từng model type
    \item Input dimension handling linh hoạt (90 features vs 900 raw inputs)
\end{itemize}

Điều này thể hiện trade-off thiết kế: complexity tăng để đổi lấy flexibility và unified user experience.

\subsection{Hạn chế TFLite và Giải pháp Keras Surrogate}
\label{sec:disc_tflite_limitations}

Finding 15 tiết lộ những hạn chế quan trọng trong hệ sinh thái ONNX-ML và TFLite đối với các thuật toán ML cổ điển.

\subsubsection{Phân tích Chi tiết về ONNX-ML Bottleneck}
Pipeline chuyển đổi bị gián đoạn tại bước ONNX$\rightarrow$TensorFlow:
\begin{enumerate}
    \item \texttt{skl2onnx}: Thành công tạo ONNX-ML operators (TreeEnsembleClassifier, SVMClassifier)
    \item \texttt{onnx2tf}: \textbf{Không hỗ trợ} ONNX-ML operators---chỉ hỗ trợ standard ONNX operators
    \item \textbf{Kết quả}: RF/SVM không thể chuyển đổi trực tiếp sang TFLite
\end{enumerate}

Đây là một hạn chế hệ thống trong hệ sinh thái, không phải lỗi triển khai cụ thể của framework.

\subsubsection{Đánh giá Keras Surrogate Method}
Giải pháp Keras surrogate đưa ra trade-offs quan trọng:

\textbf{Ưu điểm}:
\begin{itemize}
    \item Cho phép RF/SVM sử dụng TFLite pipeline
    \item Tận dụng được quantization INT8 và runtime optimizations
    \item Tương thích với hệ sinh thái TensorFlow hiện có
\end{itemize}

\textbf{Nhược điểm}:
\begin{itemize}
    \item \textbf{Approximation loss}: Keras MLP chỉ đạt 80-90\% agreement với RF/SVM gốc
    \item \textbf{Memory overhead}: MLP surrogate có thể lớn hơn direct C++ implementation
    \item \textbf{Training complexity}: Cần additional training step và hyperparameter tuning
\end{itemize}

\textbf{Khuyến nghị sử dụng}: 
\begin{itemize}
    \item \textbf{RF/SVM}: Ưu tiên direct C++ generation (exact, compact, fast)
    \item \textbf{NN/CNN}: Sử dụng TFLite để tận dụng ecosystem mature
    \item \textbf{TFLite surrogate}: Chỉ khi yêu cầu standardized runtime integration
\end{itemize}

\subsubsection{Implications cho Edge AI Ecosystem}
Findings này chỉ ra một gap quan trọng trong edge AI toolchain: ONNX-ML và standard ONNX có operator domains khác nhau, tạo ra compatibility issues ở interface layers. Điều này nhấn mạnh giá trị của direct code generation approach cho classical ML algorithms.

\section{Các Hướng Nghiên Cứu Tương Lai}

\subsection{Các Mở Rộng Ngắn Hạn}
Các cải tiến ngay lập tức trong kiến trúc hiện tại:
\begin{enumerate}
    \item \textbf{Các Thuật Toán Bổ Sung}: Tích hợp XGBoost, LightGBM, k-NN
    \item \textbf{Lựa Chọn Đặc Trưng}: Triển khai phân tích tầm quan trọng đặc trưng tự động để giảm chiều
    \item \textbf{Nén Mô Hình}: Cắt tỉa, huấn luyện nhận thức lượng tử hóa để giảm thêm bộ nhớ
    \item \textbf{Nhiều Nền Tảng Hơn}: Hỗ trợ triển khai ESP32, STM32, nRF5340
    \item \textbf{Trực Quan Hóa Thời Gian Thực}: Luồng hoạt động trực tiếp từ thiết bị đã triển khai đến GUI framework
\end{enumerate}

\subsection{Tầm Nhìn Dài Hạn}
Các cải tiến chuyển đổi yêu cầu thay đổi kiến trúc:
\begin{enumerate}
    \item \textbf{Tích Hợp Học Sâu}: Tạo TFLite Micro cho các mô hình CNN/LSTM
    \item \textbf{Học Liên Bang}: Cho phép huấn luyện mô hình cộng tác trên nhiều thiết bị mà không chia sẻ dữ liệu thô
    \item \textbf{AI Có Thể Giải Thích}: Trực quan hóa các mẫu cảm biến nào kích hoạt dự đoán hoạt động cụ thể
    \item \textbf{Hợp Nhất Đa Phương Thức}: Kết hợp IMU với âm thanh, áp suất hoặc dữ liệu camera
    \item \textbf{Học Trực Tuyến}: Các mô hình đã triển khai thích ứng với các mẫu cụ thể của người dùng theo thời gian
    \item \textbf{Tích Hợp Đám Mây}: Backend tùy chọn cho sao lưu dữ liệu, phiên bản mô hình, chia sẻ mô hình cộng đồng
\end{enumerate}

\subsection{Các Cải Tiến Khả Năng Sử Dụng}
\begin{itemize}
    \item Các nghiên cứu người dùng chính thức so sánh giao diện tiền xử lý với các lựa chọn thay thế (Edge Impulse, coding thủ công)
    \item Các quy trình làm việc có hướng dẫn cho người dùng mới với hướng dẫn từng bước và các tập dữ liệu ví dụ
    \item Giám sát thiết bị thời gian thực cho thấy các luồng cảm biến trực tiếp và kết quả dự đoán
    \item Kiểm tra phần cứng tự động để xác thực các mô hình đã triển khai trước khi triển khai hiện trường
\end{itemize}

\section{Nhận Xét Kết Thúc}

Sự phổ biến của các cảm biến đeo được và điện toán biên phổ biến tạo ra các cơ hội chưa từng có cho các ứng dụng theo dõi chuyển động trải dài y tế, thể thao, nghiên cứu và xa hơn nữa. Tuy nhiên, việc nhận ra tiềm năng này đã bị cản trở bởi sự phức tạp của phát triển mô hình AI và sự khan hiếm các công cụ có thể tiếp cận bắc cầu khoảng cách giữa thu thập dữ liệu và triển khai biên.

Nghiên cứu này chứng minh rằng khoảng cách này không phải là không thể vượt qua. Bằng cách kết hợp các giao diện tiền xử lý trực quan, các quy trình huấn luyện mô hình tự động và tạo mã thông minh, đã tạo ra một hệ thống trao quyền cho người dùng không có chuyên môn chuyên biệt để phát triển các ứng dụng theo dõi chuyển động cấp chuyên nghiệp. Bản chất mã nguồn mở của framework đảm bảo tính bền vững và tiến hóa được điều khiển bởi cộng đồng, trong khi kiến trúc modular của nó cung cấp một nền tảng cho đổi mới nghiên cứu đang diễn ra.

Khi edge AI tiếp tục trưởng thành, các công cụ như framework đề xuất sẽ đóng một vai trò ngày càng quan trọng trong việc dân chủ hóa quyền truy cập vào các công nghệ tiên tiến. Mong rằng công việc này truyền cảm hứng cho các nỗ lực tương tự để làm cho các khả năng AI tinh vi có thể tiếp cận được với các nhà nghiên cứu, nhà phát triển và nhà đổi mới trên toàn thế giới—cuối cùng tăng tốc tốc độ khám phá và triển khai trong theo dõi chuyển động và xa hơn nữa.
